\documentclass[a4paper]{report}

% Use swiss german letters
\usepackage[utf8]{inputenc}

% Language: german
\usepackage[ngerman]{babel}

% Fancy Figures
\usepackage{graphicx}

% Use Times
\usepackage{mathptmx}

% Display the Bibliography in the TOC
\usepackage{tocbibind}

% Better lists
\usepackage{enumitem}

% Use biblatex
\usepackage[style=apa,backend=biber,citestyle=authoryear]{biblatex} 

% Tell BibLatex to use the ngerman language mapping
\DeclareLanguageMapping{ngerman}{ngerman-apa}

% Define the bibliography file
\addbibresource{bibliography.bib}

% To let LaTeX handle "
\usepackage[autostyle=true, german=quotes]{csquotes}

% Titlepage
\newcommand*{\titleAP}{\begingroup % Create the command for including the title page in the document
	\centering
	\vspace*{\baselineskip} % Whitespace at the top of the page
	
	{\Large FirstName LastName}\\[0.167\textheight] % Author name
	
	{\Huge\bfseries Projektdokumentation PREN Gruppe03}\\[\baselineskip]
	
	%TODO review subtitle
	{\Large \textit{Subtitle}}\\
	\today
	
	\vspace*{3\baselineskip} % Whitespace at the bottom of the page
	\endgroup}

% Define the path were images are found
\graphicspath{{./img/}}

\begin{document}

\titleAP

\newpage

\begin{abstract}
	Hier würde man das Abstract oder Management Summary schreiben.
\end{abstract}

\tableofcontents

\newpage

\chapter{Einleitung}
\label{ch:Intro}

\section{Präambel}
Eine Einleitung in ein \LaTeX Dokument. 
Dieses ist im Moment auf Github gehosted, Git ist ein Codeverwaltungsprogramm welches von vielen Programmierern und OpenSource Projekten benutzt wird. \parencite{Git2017}

Ein leichtverständliches Buch wurde von \citeauthor{Chacon2016} geschrieben und ist frei verfügbar.

In diesem Satz verweise ich auf Bild \ref{fig:MemeMagic}, welches mir \LaTeX automatisch richtig beschriftet.

\begin{figure}[h!]
	\centering
	\includegraphics[width=0.5\linewidth,keepaspectratio]{MemeMagic}
	\caption{\LaTeX is magic - sometimes black magic...}
	\label{fig:MemeMagic}
	
\end{figure}

\chapter{Projektorganisation}

\section{Teamübersicht}

\begin{tabular}{|l|l|}
	\hline 
	\textbf{Name} & \textbf{Studium} \\ 
	\hline 
	Pascal Baumann & Informatik \\ 
	\hline 
	Basil Bachmann & Maschinenbau \\ 
	\hline 
	David Craven & Elektrotechnik \\ 
	\hline 
	Victor Guntern & Maschinenbau \\ 
	\hline
	Markus Kempf & Maschinenbau \\ 
	\hline  
	Eve Meier & Informatik \\ 
	\hline 
	Jan Odermatt & Elektrotechnik \\
	\hline
	Simon Rohrer & Maschinenbau \\
	\hline
\end{tabular} 

\section{Projektrollen}

\begin{tabular}{|l|l|l|}
	\hline 
	\textbf{Rolle} & \textbf{Aufgaben} & \textbf{Teammitglied} \\ 
	\hline 
	Projektleiter & Gesamtübersicht des Projektes halten
	Überprüfen ob Vorgaben eingehalten werden 
	Teammeetings organisieren und Informationsaustausch sicherstellen
	Kostenverwaltung & Eve \\ 
	\hline 
	Projektplaner & Aktualisieren des Terminplanes 
	Rahmenplanung und Meilensteineinhaltung 
	Pflege Taskboard & Markus \\ 
	\hline 
	Verantwortlicher Dokumentation & Meilensteinabgabe zusammenstellen
	Protokolle führen
	Unterstützung und Pflege LaTeX & Pascal \\ 
	\hline 
	Fachverantwortliche & Projektstand und Feedback
	Ansprechperson bei Fragen
	Aktualisierung Risikomanagement
	Koordination Versuche und Rechere & I Pascal, ET Jan, M Markus \\ 
	\hline 
\end{tabular} 

\section{Hilfsmittel / Tools}
\begin{tabular}{|l|l|}
	\hline 
	\textbf{Aufgabe} & \textbf{Hilfsmittel} \\ 
	\hline 
	Dokumente und Dokumentation & LaTeX / MiKTeX / GitHub \\
	\hline
	Quellen & Mendeley \\
	\hline
	Dateiablage für Teamaustausch & Dropbox \\
	\hline
	Dateiablage für Abgabe & Ilias \\
	\hline
	Projekt- und Budgetplan & MS Excel 2016 \\
	\hline
	Kommunikation Team & WhatsApp Gruppe od. HSLU Mailadresse\\
	\hline
	Aufgabenverwaltung & SCRUM-angelehntes Board, physisch, Teaminsel\\
	\hline
\end{tabular}

\section{Wochenplan}
\begin{tabular}{|l|l|}
	\hline 
	\textbf{Tag} & \textbf{Beschreibung} \\
	\hline
	DO 08:30 & Alle Mitglieder sind in der Teaminsel \\
	\hline
	DO 08:30-09:00 & Besprechung Team-intern erledigte Aufgaben, Fragen, weiteres Vorgehen \\
	\hline
	DO 09:00-10:00 & Arbeiten im Team od. selbständig \\
	\hline
	DO 10:00-10:20 & Pause \\
	\hline
	DO 10:20-12:00 & Arbeiten im Team od. selbstständig \\
	\hline
	FR 08:30 & Alle Mitglieder sind in der Teaminsel \\
	\hline
	FR 08:30-09:00 & Besprechung Team-intern, Vorbereitung Meeting mit Dozent \\
	\hline
	FR 09:00-09:30& Besprechung mit Dozent \\
	\hline
	FR 09:30-10:00 & Arbeiten im Team od. selbstständig \\
	\hline
	FR 10:00-10:20 & Pause \\
	\hline
	FR 10:20-11:00 & Arbeiten im Team od. selbstständig \\
	\hline
	FR 11:00-11:30 & Kurzbesprechung, Taskboard aktualisieren \\
	\hline
	FR 11:30-12:00 & Arbeiten im Team od. selbstständig (freiwillig)\\
	\hline
\end{tabular}
	

\chapter{Anforderungen}
\section{Projektanforderungen}
\begin{tabular}{|l|l|l|}
	\hline 
	\textbf{Nr.} & \textbf{Bezeichnung} & \textbf{Beschreibung} \\ 
	\hline 
	1.1 & Projektabgabe & Dezember 2017 \\ 
	\hline 
	1.2 & Eigenleistung & Systemkomponenten können zugekauft werden \\ 
	\hline
	1.3 & Interdisziplinarität & Disziplinen / Abteilungen arbeiten zusammen \\
	\hline
	1.4 & Lieferantenwahl & Für Sammelbestellungen gem. Kapitel 4.5 Aufgabenstellung. Wird Material vom Team selbst gekauft, können die Kosten zurückgefordert werden \\ 
	\hline 
	1.5 & Budget f. PREN & max. 500.- CHF \\ 
	\hline 
	1.6 & Teilbudget PREN1 & max. 200.- CHF \\ 
	\hline 
	1.7 & 3D-Drucker Laufzeit & max. 25 \\ 
	\hline 
	1.8 & Lasergerät Laufzeit & max. 1 \\ 
	\hline 
	1.9 & Stunden ET-Werkstattpersonal & max. 10 \\ 
	\hline 
	1.10 & Stunden M-Werkstattpersonal & max. 10 \\ 
	\hline 
	1.11 & "Gesponsorte" Komponenten & Werden mit einem realistischen Preis in die Kostenrechnung einbezogen \\
	\hline
\end{tabular} 

\section{Plattform}
\begin{tabular}{|l|l|l|}
	\hline 
	\textbf{Nr.} & \textbf{Bezeichnung} & \textbf{Beschreibung} \\
	\hline 
	2.1 & Gesamtlänge & 350 $\pm$ 2cm \\ 
	\hline 
	2.2 & Masten Abstand & Abstand zwischen den Masten 350cm $\pm$ 2cm \\
	\hline
	2.3 & Masten Masse & TODO Quadratisch? 10cm $\pm$ 1 cm \\
	\hline
	2.4 & Drahtseil & Verzinkter Stahl, Durchmesser 3mm \\ 
	\hline 
	2.5 & Seilspannung & Via Umlenkrollen durch ein Gewicht mit einer Masse von 15kg \\ 
	\hline 
	2.6 & Winkel des Seiles & TODO \\
	\hline
	2.7 & Grundplatte & Spanplatte roh oder grau gestrichen. Mit Farbresten / vorstehenden Schrauben und Nahtstellen ist zu rechnen \\
	\hline
	2.8 & Startfeld & 50cm $\pm$ 2cm, Quadratisch \\ 
	\hline 
	2.9 & Zielplatte & TODO Gesamtmass? \\ 
	\hline 
	2.10 & Zielplatte Aussehen & TODO Wie viele konzentrische Bereiche? Der innerste, helle Bereich ist quadratisch und hat eine Seitenlänge von 6cm $\pm$ 2mm. Jeder daran anschliessende konzentrische Bereich hat eine Breite von 2.5cm $\pm$ 2mm. Die Bereiche sind abwechslungsweise hell und dunkel \\ 
	\hline 
	2.11 & Zielplatte Position & Der Absetzbereich verläuft unterhalb des Seiles und ist 2cm breit. Die Zielplatte kann bis zum Startsignal verschoben werden. Befindet sich aber immer im Absetzbereich siehe Abbildung 1 der Aufgabenstellung \\
	\hline
	2.12 & Start- und Zielplatte & TODO Matt? Glänzend? \\
	\hline
	2.13 & Hindernisse & Auf der gesamten Plattform können Hindernisse stehen. Im Umkreis von mindestens 10cm um die Startposition des Ladegutes und um die Zielplatte herum sind keine Hindernisse \\
	\hline
	2.14 & Hindernisse Höhe & Die Hindernisse haben eine maximale Höhe von 20cm. \\
	\hline
\end{tabular} 

\section{Laufkatze}
\begin{tabular}{|l|l|l|}
	\hline 
	\textbf{Nr.} & \textbf{Bezeichnung} & \textbf{Beschreibung} \\
	\hline 
	3.1 & Steuerung & Autonom \\
	\hline
	3.2 & Inbetriebnahme & Darf max. 2min dauern \\
	\hline
	3.3 & Startsignal & Darf per Kopfdruck gesendet werden \\
	\hline
	3.4 & Geschwindigkeit & Um die Aufgabe zu bewältigen steht der Laufkatze ein Zeitfensters von 4min zur Verfügung. \\
	\hline
	3.5 & Aussendimensionen & Die Laufkatze darf in ihrer Projektion das Startfeld nicht überschreiten. 50cm $\pm$ 2cm x 50cm $\pm$ 2cm \\
	\hline
	3.6 & Bauart & Sämtliche Sensorik muss auf dem Gerät selbst montiert sein. \\
	\hline
	3.7 & Fahrweise & Das Gerät darf nur das Drahtseil und den zweiten Masten berühren. Die gesamte Plattform, insbesondere Drahtseil, die Last und die Zielplatte dürfen nicht beschädigt oder sonst irgendwie verändert werden. Es ist beispielsweise nicht erlaubt, Navigationshilfen anzubringen. \\
	\hline
	3.8 & Ladegut & Das Gerät muss ein Ladegut transportieren können. \\
	\hline
	3.9 & Zielerkennung & Das erkennen der Zielplatte muss selbstständig erfolgen\\
	\hline 
\end{tabular}

\section{Ladegut}
\begin{tabular}{|l|l|l|}
	\hline 
	\textbf{Nr.} & \textbf{Bezeichnung} & \textbf{Beschreibung} \\
	\hline
	4.1 & Material & Holz \\
	\hline
	4.2 & Dimensionen & Seitenlänge 5cm $\pm$ 0.5cm \\
	\hline
	4.3 & Gewicht & TODO \\
	\hline
	4.4 & Aufnahme & Metallischer, magnetischer Hacken oben in der Mitte des Würfels. Innendurchmesser des Hakens ist $\pm$ 1cm. \\
	\hline
	4.5 & Hindernisse & Das Ladegut darf Hindernisse nicht berühren. \\
	\hline
	4.6 & Position & Die Position des Ladegutes muss in Echtzeit angezeigt werden, dass der Schiedsrichter jederzeit die angezeigten Werte gut erkennen kann. TODO auf Gerät od. Extern? \\
	\hline
	4.7 & Positionsbestimmung & Die Mitte des Bodens des Ladegutes wird verwendet. Die Position muss in x- und z-Richtung bestimmt werden. Der Nullpunkt des zu verwendenden Koordinatensystems ins in Abbildung 1 der Aufgabenstellung definiert. \\
	\hline
	4.8 & Absetzen & Das Ladegut muss innerhalb des automatisch erkannten Zielbereiches abgesetzt werden. \\
	\hline
\end{tabular}

\section{Umfeld}
\begin{tabular}{|l|l|l|}
	\hline 
	\textbf{Nr.} & \textbf{Bezeichnung} & \textbf{Beschreibung} \\
	\hline
	5.1 & Licht & TODO \\
	\hline
	5.2 & Temperaturen & Bei Lagerung und Betrieb Zimmertemperatur 15C - 20C TODO Grad Zeichen\\
	\hline
\end{tabular}

\chapter{Risikomanagement}
In diesem Kapitel werden mögliche Risiken während des Projektverlaufes aufgelistet. Dabei werden Projektrisiken nummeriert. Ihre Eintrittswahrscheinlichkeit und ihr Schadensausmass wird eingeschätzt. Besteht ein grosses Risiko, werden zusätzlich Massnahmen definiert. 
\section{Definitionen}
TODO Eintrittswahrscheinlichkeit und Schadensausmass definieren. 
1-4 ?
\section{Risikokatalog}
\begin{tabular}{|l|l|l|l|}
	\hline 
	\textbf{Nr.} & \textbf{Risiko} & \textbf{Wahrscheinlichkeit} & \textbf{Schadensausmass} \\
	\hline
\end{tabular}

\section{Risikobewertung}
TODO tolles farbiges Matrixteil yeah
 
\section{Massnahmen}
TODO Bei Risiken im Roten und Gelben Bereich Massnahmen definieren 

\newpage

\printbibliography

\end{document}